\documentclass[12pt]{article}
\usepackage{amsmath, amssymb, amsthm}
\usepackage{geometry}
\usepackage{booktabs}
\usepackage{hyperref}
\usepackage{natbib}
\geometry{margin=1in}

\newtheorem{definition}{Definition}
\newtheorem{remark}{Remark}
\newtheorem{example}{Example}

\title{Shapley Values, Causal Shapley Values, and Responsibility Attribution:\\
A Critical Analysis}
\date{}

\begin{document}

\maketitle

\tableofcontents
\newpage

%% ============================================================
\section{Introduction}
%% ============================================================

Shapley values have become one of the most widely used tools for explaining the predictions
of complex machine learning models. Grounded in cooperative game theory, they provide a
principled decomposition of a model's prediction into additive contributions from individual
features. However, when the goal is \emph{causal responsibility attribution} --- explaining
not just what a model predicts but why an outcome occurred --- standard Shapley values
suffer from fundamental limitations that are not always made explicit in the literature.

This chapter develops these limitations carefully, starting from first principles. We work
throughout with a simple but instructive running example: a linear Gaussian chain
$X \to M \to Y$, where all variables are Gaussian and the relationships are additive.
This model is simple enough to permit exact analytic computation while being rich enough
to expose the core conceptual issues. We then examine how Causal Shapley values
\citep{heskes2020causal} partially address these issues, where they succeed, and where
fundamental gaps remain.

%% ============================================================
\section{The Running Model}
%% ============================================================

Throughout this chapter we work with the following generative model:

\begin{align}
X &\sim \mathcal{N}(0.5,\ 0.25) \label{eq:X}\\
M &= X + \varepsilon_M, \quad \varepsilon_M \sim \mathcal{N}(0,\ 0.1) \label{eq:M}\\
Y &= M + \varepsilon_Y, \quad \varepsilon_Y \sim \mathcal{N}(0,\ 0.1) \label{eq:Y}
\end{align}

The causal graph is $X \to M \to Y$, a simple mediation chain. All noise terms are
independent of each other and of $X$. Since everything is jointly Gaussian, all conditional
expectations are exact linear functions --- no approximations are required.

Key distributional facts:
\begin{align}
\mathbb{E}[X] = \mathbb{E}[M] = \mathbb{E}[Y] &= 0.5 \\
\text{Var}(X) = 0.25, \quad \text{Var}(M) &= 0.35, \quad \text{Var}(Y) = 0.45 \\
\text{Cov}(X,M) = 0.25, \quad \text{Cov}(X,Y) &= 0.25, \quad \text{Cov}(M,Y) = 0.35
\end{align}

\subsection{Prediction Functions}

For each choice of target variable, the optimal predictor (minimizing squared error) is the
conditional expectation. We derive these exactly.

\paragraph{Target $Y$, features $\{X, M\}$:}
Since $Y = M + \varepsilon_Y$ with $\varepsilon_Y$ independent of $X$ and $M$:
\begin{equation}
f_Y(X, M) = \mathbb{E}[Y \mid X, M] = M
\end{equation}
$X$ carries no additional information about $Y$ once $M$ is known --- $M$ screens $X$ from
$Y$ by d-separation.

\paragraph{Target $M$, features $\{X, Y\}$:}
Using the multivariate Gaussian conditional formula with the covariance structure above:
\begin{equation}
f_M(X, Y) = \mathbb{E}[M \mid X, Y] = 0.5X + 0.5Y
\end{equation}
$M$ is exactly halfway between $X$ and $Y$ because the noise variances at each step are equal,
making $X$ and $Y$ equally informative about $M$.

\paragraph{Target $X$, features $\{M, Y\}$:}
Since $Y = M + \varepsilon_Y$, conditioning on $M$ renders $Y$ independent of $X$
(d-separation: $M$ blocks the only path $X \to M \to Y$). Therefore:
\begin{equation}
f_X(M, Y) = \mathbb{E}[X \mid M, Y] = \mathbb{E}[X \mid M]
= 0.5 + \frac{5}{7}(M - 0.5)
\end{equation}
$Y$ drops out exactly. The coefficient $\frac{5}{7} = \frac{\text{Cov}(X,M)}{\text{Var}(M)}
= \frac{0.25}{0.35}$ reflects that $M$ is a noisy signal of $X$: observing $M=1$ raises our
estimate of $X$ to $0.857$ but not all the way to $1$, because some of $M$'s value is due
to $\varepsilon_M$.

%% ============================================================
\section{Shapley Values: Definition and Properties}
%% ============================================================

\subsection{Cooperative Game Theory Background}

A cooperative game is a pair $(N, v)$ where $N = \{1, \ldots, n\}$ is the set of players
and $v: 2^N \to \mathbb{R}$ is the \emph{characteristic function} mapping each subset
(coalition) $S \subseteq N$ to the value that coalition can produce. The Shapley value
\citep{shapley1953value} is the unique attribution satisfying four axioms:
\begin{enumerate}
\item \textbf{Efficiency}: $\sum_i \phi_i = v(N) - v(\emptyset)$
\item \textbf{Symmetry}: If $v(S \cup \{i\}) = v(S \cup \{j\})$ for all $S$, then
      $\phi_i = \phi_j$
\item \textbf{Dummy}: If $v(S \cup \{i\}) = v(S)$ for all $S$, then $\phi_i = 0$
\item \textbf{Linearity}: For two games $v_1, v_2$: $\phi_i(\alpha_1 v_1 + \alpha_2 v_2)
      = \alpha_1 \phi_i(v_1) + \alpha_2 \phi_i(v_2)$
\end{enumerate}

The unique solution is:
\begin{equation}
\phi_i = \sum_{S \subseteq N \setminus \{i\}}
\frac{|S|!\,(|N|-|S|-1)!}{|N|!}
\bigl[v(S \cup \{i\}) - v(S)\bigr]
\label{eq:shapley}
\end{equation}

The weight $\frac{|S|!(|N|-|S|-1)!}{|N|!}$ equals the probability that, in a uniformly
random ordering of all players, exactly the members of $S$ precede $i$. Equivalently,
$\phi_i$ is the average marginal contribution of $i$ across all orderings.

\subsection{SHAP: Shapley Values for Machine Learning}

\citet{lundberg2017unified} apply Shapley values to explain model predictions by defining:
\begin{itemize}
\item Players = input features
\item $v(S) = \mathbb{E}_{X_{\bar{S}}}\bigl[f(X_S = x_S,\, X_{\bar{S}})\bigr]$
\end{itemize}
where $\bar{S} = N \setminus S$ and missing features are marginalized over their
(marginal or conditional) distribution. The attribution decomposes:
\begin{equation}
f(x) = \underbrace{\mathbb{E}[f(X)]}_{f_0} + \sum_{i=1}^n \phi_i
\end{equation}

\begin{remark}[SHAP explains predictions, not outcomes]
SHAP is defined via the \emph{local accuracy} property: the explanation model must match
$f(x)$, the model prediction, not the actual outcome $Y$. If $f(x) = \mathbb{E}[Y \mid X=x]$
(as in regression), then SHAP explains $\mathbb{E}[Y \mid X=x] - \mathbb{E}[Y]$. The
residual $Y - f(x)$ --- the noise --- is entirely invisible to SHAP by construction. This
is not stated as a limitation in \citet{lundberg2017unified}; it follows from the definition.
\end{remark}

%% ============================================================
\section{Plain SHAP on the Running Model}
%% ============================================================

We compute SHAP for all three targets at the instance $X=1, M=1, Y=1$, with baseline
$\mathbb{E}[\cdot] = 0.5$ throughout. The value function uses observational conditioning:
$v(S) = \mathbb{E}[f \mid X_S = x_S]$.

\subsection{Target $Y$, Features $\{X, M\}$}

\begin{align}
v(\emptyset) &= \mathbb{E}[M] = 0.5 \\
v(\{X\}) &= \mathbb{E}[M \mid X=1] = 1 \\
v(\{M\}) &= \mathbb{E}[M \mid M=1] = 1 \\
v(\{X,M\}) &= f_Y(1,1) = 1
\end{align}

Applying equation~\eqref{eq:shapley} with $|N|=2$:
\begin{equation}
\phi_X = \tfrac{1}{2}(1 - 0.5) + \tfrac{1}{2}(1 - 1) = 0.25
\qquad
\phi_M = \tfrac{1}{2}(1 - 0.5) + \tfrac{1}{2}(1 - 1) = 0.25
\end{equation}

\subsection{Target $M$, Features $\{X, Y\}$}

\begin{align}
v(\emptyset) &= 0.5 \\
v(\{X\}) &= \mathbb{E}[0.5X + 0.5Y \mid X=1] = 0.5 + 0.5\cdot\mathbb{E}[Y\mid X=1] = 1 \\
v(\{Y\}) &= \mathbb{E}[0.5X + 0.5Y \mid Y=1] = 0.5\cdot\mathbb{E}[X\mid Y=1] + 0.5
         = 0.5\cdot 0.778 + 0.5 = 0.889 \\
v(\{X,Y\}) &= 1
\end{align}

where $\mathbb{E}[X \mid Y=1] = 0.5 + \frac{0.25}{0.45}(0.5) = 0.778$.

\begin{equation}
\phi_X = \tfrac{1}{2}(1-0.5) + \tfrac{1}{2}(1-0.889) = 0.306
\qquad
\phi_Y = \tfrac{1}{2}(0.889-0.5) + \tfrac{1}{2}(1-1) = 0.194
\end{equation}

\subsection{Target $X$, Features $\{M, Y\}$}

\begin{align}
v(\emptyset) &= 0.5 \\
v(\{M\}) &= 0.5 + \tfrac{5}{7}(0.5) = 0.857 \\
v(\{Y\}) &= \mathbb{E}[X \mid Y=1] = 0.5 + \tfrac{0.25}{0.45}(0.5) = 0.778 \\
v(\{M,Y\}) &= \mathbb{E}[X \mid M=1, Y=1] = 0.857 \quad \text{(since $X \perp Y \mid M$)}
\end{align}

\begin{equation}
\phi_M = \tfrac{1}{2}(0.857-0.5) + \tfrac{1}{2}(0.857-0.778) = 0.219
\qquad
\phi_Y = \tfrac{1}{2}(0.778-0.5) + \tfrac{1}{2}(0.857-0.857) = 0.139
\end{equation}

\subsection{Summary: Plain SHAP}

\begin{table}[h]
\centering
\begin{tabular}{llccc}
\toprule
Target & Method & $\phi_X$ & $\phi_M$ & $\phi_Y$ \\
\midrule
$Y$ & Plain SHAP & 0.25  & 0.25  & ---   \\
$M$ & Plain SHAP & 0.306 & ---   & 0.194 \\
$X$ & Plain SHAP & ---   & 0.219 & 0.139 \\
\bottomrule
\end{tabular}
\caption{Plain SHAP attributions for instance $X=1, M=1, Y=1$.}
\end{table}

\begin{remark}[Causal direction is invisible to plain SHAP]
For target $X$, plain SHAP assigns $\phi_Y = 0.139 > 0$, even though $Y$ does not cause $X$.
This arises because $v(\{Y\}) = 0.778 > 0.5$: observing $Y=1$ is statistically informative
about $X=1$ through the chain $X \to M \to Y$. SHAP cannot distinguish ``$Y$ is informative
about $X$'' from ``$Y$ causes $X$'' because the value function is grounded in observational
conditionals, which are symmetric under reversal of causal direction.
\end{remark}

%% ============================================================
\section{The Mediator Dilution Problem}
%% ============================================================

A second fundamental limitation of plain SHAP concerns the \emph{stability of attribution
under model reparametrization}.

Consider expanding the chain $X \to M \to Y$ to $X \to M_1 \to M_2 \to Y$ by splitting
the single intermediary $M$ into two steps, each with half the noise variance:
\begin{align}
M_1 &= X + \varepsilon_1, \quad \varepsilon_1 \sim \mathcal{N}(0, 0.05) \\
M_2 &= M_1 + \varepsilon_2, \quad \varepsilon_2 \sim \mathcal{N}(0, 0.05) \\
Y &= M_2 + \varepsilon_Y, \quad \varepsilon_Y \sim \mathcal{N}(0, 0.1)
\end{align}

The total noise $\text{Var}(Y \mid X) = 0.05 + 0.05 + 0.1 = 0.2$ is identical to the
original model. The joint distribution of $(X, Y)$ is unchanged.

Since $f_Y(X, M_1, M_2) = M_2$, each variable individually predicts $Y$ almost perfectly
(low noise). As a result, SHAP splits the fixed attribution budget $f(x) - f_0 = 0.5$
equally across all three features:

\begin{equation}
\phi_X = \phi_{M_1} = \phi_{M_2} = \tfrac{0.5}{3} \approx 0.167
\end{equation}

More generally, with $k$ intermediaries: $\phi_X = \frac{0.5}{k+1}$.

\begin{table}[h]
\centering
\begin{tabular}{lcc}
\toprule
Model & $\phi_X$ & Interpretation \\
\midrule
$X \to M \to Y$              & 0.25  & $X$ gets 50\% of total attribution \\
$X \to M_1 \to M_2 \to Y$   & 0.167 & $X$ gets 33\% \\
$X \to M_1 \to M_2 \to M_3 \to Y$ & 0.125 & $X$ gets 25\% \\
\bottomrule
\end{tabular}
\caption{SHAP attribution to $X$ as a function of chain length, with identical $(X,Y)$
joint distribution throughout.}
\end{table}

\begin{remark}[Mediator dilution]
The attribution to $X$ decreases monotonically as the causal chain is refined, even though
the causal relationship between $X$ and $Y$ is unchanged. This is because each additional
mediator enters as a new player in the Shapley game and claims a share of the fixed budget.
The attribution is determined by modeling granularity, not by facts about causal responsibility.
\end{remark}

We note that this problem can be given a natural illustration in terms of a \emph{telephone
game}: $X$ is the original message, $M_1, M_2, \ldots$ are successive intermediaries each
introducing small distortions, and $Y$ is the final received message. Splitting one
intermediary into two --- while keeping the total distortion constant --- reduces the
apparent responsibility of the original sender under SHAP.

%% ============================================================
\section{Causal Shapley Values}
%% ============================================================

\subsection{Definition}

\citet{heskes2020causal} propose replacing the observational conditional in the value
function with an interventional one using Pearl's do-calculus:

\begin{equation}
v_{\text{causal}}(S) = \mathbb{E}\bigl[f(X) \mid do(X_S = x_S)\bigr]
= \int P(X_{\bar{S}} \mid do(X_S = x_S))\, f(X_{\bar{S}}, x_S)\, dX_{\bar{S}}
\label{eq:causal_v}
\end{equation}

By the rules of do-calculus, intervening on $X_S$ only affects the distribution of
descendants of $X_S$. For non-descendants, $P(X_j \mid do(X_S = x_S)) = P(X_j)$.
Therefore equation~\eqref{eq:causal_v} is equivalent to:

\begin{equation}
v_{\text{causal}}(S) = \mathbb{E}_{X_{\bar{S},\,\text{non-desc}} \sim P(\cdot),\;
X_{\bar{S},\,\text{desc}} \sim P(\cdot \mid do(X_S=x_S))}\bigl[f(X_{\bar{S}}, x_S)\bigr]
\end{equation}

The Shapley formula~\eqref{eq:shapley} is otherwise unchanged. The method retains all four
Shapley axioms.

\subsection{Causal Shapley on the Running Model}

In the model $X \to M \to Y$ with no confounding, the interventional and observational
conditionals coincide for $X$ (since $X$ has no parents). However they differ for $M$ and $Y$
when these appear as out-of-coalition features with ancestors in the coalition.

\paragraph{Target $Y$, features $\{X, M\}$:}

\begin{align}
v(\{X\}) &= \mathbb{E}_{M \sim P(M \mid do(X=1))}[M] = \mathbb{E}[M \mid do(X=1)] = 1
\end{align}

With no confounding this equals the plain SHAP value. Result: $\phi_X = \phi_M = 0.25$.

\paragraph{Target $M$, features $\{X, Y\}$:}

$Y$ is a descendant of $X$, so when $X \in S$:
\begin{align}
v(\{X\}) &= \mathbb{E}_{Y \sim P(Y \mid do(X=1))}[0.5\cdot 1 + 0.5Y]
         = 0.5 + 0.5 \cdot 1 = 1
\end{align}

When $Y \in S$ without $X$: $X$ is not a descendant of $Y$, so $X$ is marginalized over $P(X)$:
\begin{align}
v(\{Y\}) &= \mathbb{E}_{X \sim P(X)}[0.5X + 0.5\cdot 1] = 0.5\cdot 0.5 + 0.5 = 0.75
\end{align}

This differs from plain SHAP ($v(\{Y\}) = 0.889$). Result:
\begin{equation}
\phi_X = 0.375, \qquad \phi_Y = 0.125
\end{equation}

\paragraph{Target $X$, features $\{M, Y\}$:}

$Y$ is a descendant of $M$, but $f_X$ does not depend on $Y$. Moreover $do(Y=y)$ cuts the
$M \to Y$ edge and does not affect $M$'s distribution. So $v(\{Y\}) = 0.5 = v(\emptyset)$.
Result:
\begin{equation}
\phi_M = 0.357, \qquad \phi_Y = 0
\end{equation}

\begin{table}[h]
\centering
\begin{tabular}{llccc}
\toprule
Target & Method & $\phi_X$ & $\phi_M$ & $\phi_Y$ \\
\midrule
$Y$ & Plain SHAP       & 0.25  & 0.25  & ---   \\
$Y$ & Causal SHAP      & 0.25  & 0.25  & ---   \\
$M$ & Plain SHAP       & 0.306 & ---   & 0.194 \\
$M$ & Causal SHAP      & 0.375 & ---   & 0.125 \\
$X$ & Plain SHAP       & ---   & 0.219 & 0.139 \\
$X$ & Causal SHAP      & ---   & 0.357 & 0     \\
\bottomrule
\end{tabular}
\caption{Comparison of plain and causal SHAP attributions for instance $X=1, M=1, Y=1$.}
\end{table}

Causal Shapley correctly zeros out $\phi_Y$ for target $X$: since $do(Y=y)$ does not affect
$M$, and $f_X$ does not depend on $Y$, $Y$ adds nothing to any coalition under intervention.
This is causally correct. However $\phi_Y = 0.125 > 0$ persists for target $M$ --- a residual
we examine in the next section.

%% ============================================================
\section{Direct and Indirect Effect Decomposition}
%% ============================================================

\citet{heskes2020causal} further decompose each $\phi_i(\pi)$ into direct and indirect
components. For permutation $\pi$ with $S = \{j : j \prec_\pi i\}$:

\begin{align}
\phi_i(\pi) &= \underbrace{\mathbb{E}[f(X_{\bar{S}}, x_{S\cup i}) \mid do(X_S=x_S)]
              - \mathbb{E}[f(X_{\bar{S}\cup i}, x_S) \mid do(X_S=x_S)]}_{\text{direct effect}}
\notag\\
&\quad + \underbrace{\mathbb{E}[f(X_{\bar{S}}, x_{S\cup i}) \mid do(X_{S\cup i}=x_{S\cup i})]
              - \mathbb{E}[f(X_{\bar{S}}, x_{S\cup i}) \mid do(X_S=x_S)]}_{\text{indirect effect}}
\label{eq:decomp}
\end{align}

The \emph{direct effect} measures the change in prediction from fixing $X_i = x_i$ while
holding the distribution of other out-of-coalition features constant. The \emph{indirect
effect} measures the additional change in the distribution of out-of-coalition descendants
caused by the intervention $do(X_i = x_i)$.

\begin{remark}
``Direct'' and ``indirect'' here refer to information flow through the prediction function
$f$, not to direct edges in the causal graph. A feature with no outgoing edges in the DAG
can still have a nonzero indirect effect in this decomposition if $f$ depends on its
descendants.
\end{remark}

Applying decomposition~\eqref{eq:decomp} to our model:

\begin{table}[h]
\centering
\begin{tabular}{llccc}
\toprule
Target & Feature & Direct & Indirect & Total \\
\midrule
$Y$ & $X$ & 0      & 0.25  & 0.25  \\
$Y$ & $M$ & 0.25   & 0     & 0.25  \\
$M$ & $X$ & 0.25   & 0.125 & 0.375 \\
$M$ & $Y$ & 0.125  & 0     & 0.125 \\
$X$ & $M$ & 0.357  & 0     & 0.357 \\
$X$ & $Y$ & 0      & 0     & 0     \\
\bottomrule
\end{tabular}
\caption{Direct and indirect decomposition of causal Shapley values.}
\end{table}

Key observations:
\begin{enumerate}
\item $X$'s contribution to $Y$ is \emph{entirely indirect} (0.25 indirect, 0 direct),
      correctly reflecting that $X \to Y$ has no direct edge --- $X$ only reaches $Y$ via $M$.
\item $M$'s contribution to $Y$ is \emph{entirely direct} (0.25 direct, 0 indirect),
      correctly reflecting the direct edge $M \to Y$ with no further descendants.
\item $Y$'s contribution to $M$ (0.125) is \emph{entirely direct} and has no causal
      interpretation: $Y$ has no descendants so it cannot have indirect effects. This
      attribution arises purely because $f_M(X,Y) = 0.5X + 0.5Y$ uses $Y$ as a
      statistical signal about $M$. It reflects epistemic dependence, not causal influence.
\end{enumerate}

%% ============================================================
\section{Residual Limitations for Responsibility Attribution}
%% ============================================================

\subsection{The Epistemic vs.\ Causal Direct Effect Problem}

The persistence of $\phi_Y = 0.125$ for target $M$ under causal SHAP reveals a limitation
that the direct/indirect decomposition makes precise. The ``direct'' effect of $Y$ on the
prediction of $M$ is real --- $f_M$ genuinely depends on $Y$ as a feature --- but it is
not causal. It arises because $Y$ is statistically informative about $M$: observing $Y=1$
updates our estimate of $M$ upward even given $X$, because $Y = M + \varepsilon_Y$ carries
signal about $M$.

Causal Shapley eliminates spurious \emph{indirect} effects in the wrong causal direction
by using $do(\cdot)$ in the marginalization. But it cannot eliminate spurious \emph{direct}
effects, because those are determined by $f$ itself --- which is derived from the observational
joint distribution and therefore encodes statistical dependencies regardless of their
causal direction.

\subsection{SHAP Explains Predictions, Not Outcomes}

Consider the instance $X = 0.5, M = 0.5, Y = 1$. Here $Y = 1$ despite both $X$ and $M$
being at baseline ($0.5$), because $\varepsilon_Y = 0.5$.

The model prediction is $f_Y(0.5, 0.5) = 0.5 = f_0$. Therefore:
\begin{equation}
\phi_X = \phi_M = 0, \qquad \text{(both plain and causal SHAP)}
\end{equation}

The gap between the actual outcome ($Y=1$) and the baseline ($0.5$) is $0.5$, entirely
attributable to $\varepsilon_Y$. SHAP explains $f(x) - f_0 = 0$, leaving the full gap
unexplained. This is not a failure of SHAP --- it correctly explains the prediction ---
but it reveals that SHAP provides \emph{no information} about responsibility for the
actual realized outcome when that outcome deviates from its prediction due to noise.

\subsection{The Responsibility of Noise}

More subtly: even when features take values that differ from baseline, the decomposition
$Y - f_0 = [f(x) - f_0] + [Y - f(x)]$ splits the outcome into an explained part (which
SHAP distributes among features) and a residual (which SHAP ignores). For causal responsibility
attribution, the residual is not noise to be discarded --- it is the part of the outcome
due to unmodeled or stochastic mechanisms, and assigning it to ``nobody'' is itself a
responsibility claim.

\subsection{The Choice of Responsibility Principle}

A deeper issue is that Shapley values, even causal ones, implement a specific implicit
principle of responsibility: \emph{symmetric informational contribution, corrected for
causal structure in the marginalization}. This does not correspond to any of the standard
philosophical principles of causal responsibility:

\begin{itemize}
\item \textbf{Exogeneity / causal primacy}: assign responsibility to root causes ($X$),
      treat mediators as mechanisms
\item \textbf{Proximate causation}: assign responsibility to the most immediate cause ($M$
      for $Y$)
\item \textbf{Path-specific responsibility}: decompose along causal paths, not features
\item \textbf{Agency / manipulability}: assign responsibility to variables under an agent's
      control
\item \textbf{Actual causation}: responsibility tracks counterfactual dependence in the
      actual world
\end{itemize}

Each principle leads to different attributions. Causal SHAP does not commit to any of them
explicitly. The equal split $\phi_X = \phi_M = 0.25$ for target $Y$ is defensible under a
symmetric informational principle but is wrong under exogeneity (where $X$ should receive
all credit) and wrong under proximate causation (where $M$ should receive all credit).

%% ============================================================
\section{Conclusion}
%% ============================================================

We have worked through the theory and computation of plain and causal Shapley values on a
simple linear Gaussian chain, and identified several fundamental limitations for causal
responsibility attribution:

\begin{enumerate}
\item \textbf{Causal direction is invisible to plain SHAP}: the observational value function
      cannot distinguish $X \to Y$ from $Y \to X$.
\item \textbf{Mediator dilution}: adding intermediate variables to a causal chain reduces
      the root cause's apparent attribution without any change to the underlying causal
      relationship.
\item \textbf{Causal SHAP partially corrects direction}: by replacing conditioning with
      intervention in the marginalization, causal SHAP correctly zeros out downstream
      variables as causes of upstream ones in many cases.
\item \textbf{Epistemic direct effects persist}: causal SHAP cannot remove the attribution
      to downstream variables that appear as features in $f$, because those attributions
      arise from the prediction function itself, not from the marginalization.
\item \textbf{SHAP explains predictions, not outcomes}: by the local accuracy axiom, SHAP
      is defined to explain $f(x) - f_0$. The gap between actual outcomes and predictions
      --- the noise --- is invisible by construction.
\item \textbf{No responsibility principle is made explicit}: Shapley values implement
      symmetric informational attribution, which does not correspond to standard causal
      responsibility principles. The choice of principle must be made external to the SHAP
      framework.
\end{enumerate}

These limitations are not arguments against Shapley values as explanation tools for
predictive models --- for that purpose they are principled and well-axiomatized. They are
arguments that the step from ``explaining a prediction'' to ``attributing causal
responsibility'' requires additional theoretical machinery that Shapley values alone do not
provide.

\bibliographystyle{plainnat}
\begin{thebibliography}{9}

\bibitem[Heskes et al.(2020)]{heskes2020causal}
Heskes, T., Sijben, E., Bucur, I.G., and Claassen, T. (2020).
\newblock Causal Shapley values: Exploiting causal knowledge to explain
individual predictions of complex models.
\newblock \emph{Advances in Neural Information Processing Systems}, 33.

\bibitem[Lundberg and Lee(2017)]{lundberg2017unified}
Lundberg, S.M. and Lee, S.-I. (2017).
\newblock A unified approach to interpreting model predictions.
\newblock \emph{Advances in Neural Information Processing Systems}, 30.

\bibitem[Pearl(2009)]{pearl2009causality}
Pearl, J. (2009).
\newblock \emph{Causality: Models, Reasoning, and Inference}.
\newblock Cambridge University Press.

\bibitem[Shapley(1953)]{shapley1953value}
Shapley, L.S. (1953).
\newblock A value for n-person games.
\newblock \emph{Contributions to the Theory of Games}, 2:307--317.

\end{thebibliography}

\end{document}
